% Genealogía operatoria añadida en REV.2.
% Residencia única del principio de contracción de fibras y memoria.

\chapter[Arithmetic, discrete calculus and genealogical memory]{Operator
genealogy of arithmetic and discrete calculus: composition, contraction
of fibers, and lifts with memory}
\label{chap:genealogia-operaciones-rev2}
\index{operation!genealogy}
\index{memory!operatorial lift}

Arithmetic publishes results; a genealogical theory must additionally preserve
the conditions that made each result possible. This distinction
does not add a new operation to sum, product, power, root,
logarithm, differentiation, or integration. It specifies what structure
each contracts and what data a lift must transport in order to distinguish
antecedents identified by the same output. APP exhibits the first case on
a common support; TRIT locally orients the transition; TPK preserves
provenance throughout composition.

\section{Published operation and genealogical lifting}

Let \(X\) and \(Y\) be sets and let \(f:X\to Y\) be a map that publishes
an operation. An \emph{operation memory} is a set \(M_f\) together
with a map

\begin{equation}
 \widetilde f:X\longrightarrow Y\times M_f,
 \qquad
 \operatorname{pr}_Y\circ\widetilde f=f.
 \label{eq:levantamiento-genealogico-operacion-rev2}
\end{equation}

The first component publishes the output; the second records the data that the
reader elects to preserve. The lift is \emph{complete on the fibers}
if, for every \(y\in Y\), the restriction

\[
 \operatorname{pr}_{M_f}\circ\widetilde f\big|_{f^{-1}(y)}
 :f^{-1}(y)\longrightarrow M_f
\]

is injective. This condition expresses exactly that two antecedents with the same visible output remain distinguishable by their memory. In HMT realizations, no single token is required to preserve the entire history: the complete memory may be distributed among sheet, quotient, carry, phase, orientation, boundary, and the enriched chronological register.

If \(f:X\to Y\) and \(g:Y\to Z\) possess fixed lifts, their composition
admits the canonical lift associated with those lifts,

\begin{equation}
 \widetilde{g\circ f}:X\longrightarrow Z\times M_f\times M_g,
 \qquad
 \widetilde{g\circ f}(x)
 =\bigl(g(f(x)),m_f(x),m_g(f(x))\bigr),
 \label{eq:composicion-memorias-operacion-rev2}
\end{equation}

where \(m_f\) and \(m_g\) are the memory components. Concatenation does not
erase the provenance of the first stage in order to replace it with the second: it
coordinates them. This is the operator principle that the TPK realizes on paths.

\begin{lemma}[Composition of complete surveys]
\label{lem:composicion-levantamientos-completos-rev2}
If \(\widetilde f\) and \(\widetilde g\) are complete over the fibers of
\(f\) and \(g\), respectively, then
\(\widetilde{g\circ f}\) is complete over the fibers of \(g\circ f\).
\end{lemma}

\begin{proof}
Let \(x_1,x_2\in X\) with
\(g(f(x_1))=g(f(x_2))\) and with equal memory components in
\eqref{eq:composicion-memorias-operacion-rev2}. Equality of the memories
\(m_g(f(x_1))=m_g(f(x_2))\), together with the completeness of
\(\widetilde g\) on that fiber of \(g\), implies
\(f(x_1)=f(x_2)\). The equality
\(m_f(x_1)=m_f(x_2)\) and the completeness of \(\widetilde f\) then imply
\(x_1=x_2\).
\end{proof}

\section{Sum and product: alternatives and composition}

In a path algebra, the sum combines alternatives with compatible endpoints,

\[
 a_1\gamma_1+a_2\gamma_2,
\]

while the product concatenates stages when the end of the first
coincides with the origin of the second,

\[
 \gamma_2\gamma_1,
 \qquad t(\gamma_1)=s(\gamma_2).
\]

Addition preserves the linear coexistence of histories; multiplication constructs a composite history and, when only its value is published, may forget the factorization that produced it. On the APP chart, the evaluations \(\Sigma\) and \(\Pi\) act on the same state \((i,j)\) but organize distinct fibers. In the central block \(B_c\) of APP, the spectral separation developed in the following chapter measures precisely that local incompatibility; it is not a discrepancy between two independent tables, nor is it presented as a global distance among all readers.

On the additive sheet, translation \(T_a(x)=x+a\) is reversible on the residual
group. On the multiplicative sheet, \(M_a(x)=ax\) is invertible in
\(\mathbb Z/9\mathbb Z\) exactly when \(\gcd(a,9)=1\). The nonunitary
classes \(\overline0,\overline3,\overline6\), published in the positive
chart as \(9,3,6\), fold fibers and make visible the need to
preserve quotient and provenance. Sum and product thus constitute the two
inaugural readers of a single path geometry.

\begin{lemma}[Affine conjugation in the unit sector]
\label{lem:conjugacion-afin-app-rev3}
Let \(R=\mathbb Z/9\mathbb Z\). For \(a\in R^\times\) and \(b\in R\), the
maps \(T_b(x)=x+b\) and \(M_a(x)=ax\) satisfy
\begin{equation}
 M_aT_bM_a^{-1}=T_{ab}.
 \label{eq:conjugacion-afin-app-rev3}
\end{equation}
If \(a\in(3)\), \(M_a\) ceases to be invertible and conjugation is not
defined: its image remains contained in the radical and several branches of the
additive sheet are identified.
\end{lemma}

\begin{proof}
For \(a\in R^\times\),
\(
 M_aT_bM_a^{-1}(x)=a(a^{-1}x+b)=x+ab=T_{ab}(x)
\).
If \(a\in(3)\), the kernel of \(M_a\) is nontrivial; hence there is no
\(M_a^{-1}\), and the conjugation expression is ill-typed. The case
\(a=3\) explicitly realizes the folding through
\(\ker M_3=\operatorname{im}M_3=(3)\).
\end{proof}

\section{Power, logarithm, and root}

Let \(x\) be an element of a monoid, or of an associative algebra, and let
\(n\geq1\). The power records iteration of a composition:

\[
 x^n=\underbrace{x\cdots x}_{n\ \mathrm{factores}}.
\]

The exponent is the depth of that iteration, and the orbit of \(x\) determines
which outputs are identified. In the positive real domain, if \(b>0\),
\(b\neq1\), and \(x>0\), then
\[
 \log_b(x^n)=n\log_b x.
\]
In the complex domain, the same identity requires a compatible branch and
requires the iteration not to cross its cut. The logarithm coordinates the depth
\(n\) only after the generator \(x\) has also been fixed and
\(\log_b x\neq0\) has been required. Base, generator, domain, branch, and cut form part of the
reader.

The root inverts a power as a fibered problem,

\[
 y^n=x.
\]

The fiber may be empty, a singleton, or multiple. Selecting a root requires branch, orientation, phase, or boundary data. The subsequent construction of an internal operator \(J_E\in\operatorname{End}(V)\), with \(J_E^2=-I_V\), exemplifies this principle: the symbol \(\sqrt{-1}\) publishes a coordinate, whereas the operator, its domain, and its action preserve the structure that supports it.

\section{Discrete Calculus on Compatible Cylinders: Differentiation and Integration Operators}

Let \(P_n\) be a connected oriented path and let \(A\) be an abelian group. In the
discrete measure complex, the combinatorial coboundary

\[
 d:C^0(P_n;A)\longrightarrow C^1(P_n;A),
 \qquad (df)(e)=f(t(e))-f(s(e)),
\]

has as its kernel exactly the constant \(0\)-cochains. Reconstructing
\(f\) from \(df\) requires a boundary condition, for example, the value at a
reference vertex.

The coboundary depends on incidence and orientation, but not on a
length. If \(A\) is also a real vector space and
\(\ell:E(P_n)\to\mathbb R_{>0}\) is an edge length, the metric
derivative is defined by
\[
 D_\ell f(e)=\frac{(df)(e)}{\ell(e)}.
\]
For a route \(\gamma=(e_1,\ldots,e_r)\) compatible with the orientations,
\begin{equation}
 \sum_{k=1}^r D_\ell f(e_k)\,\ell(e_k)
 =f(t(\gamma))-f(s(\gamma)).
 \label{eq:integracion-derivada-metrica-rev2}
\end{equation}
Integration thus accumulates a cochain along a route and recovers the
endpoint difference when the cochain is exact. Coboundary, metric derivative,
and integral depend on distinct data: incidence and orientation; length and
scale; path and boundary section.

This dependence explains why the passage from counting marks to measuring intervals
is structural. The derivation lives on interstitial relations;
integration composes those relations; the cycle adds the topological carrier of
holonomy; TPK transport determines which memory survives the revolution.

\section{Hierarchy of fibers and function of the TPK}

The hierarchy can be condensed without identifying its levels:

\begin{center}
\small
\begin{tabularx}{\textwidth}{@{}
  >{\raggedright\arraybackslash}p{0.13\textwidth}
  >{\raggedright\arraybackslash}p{0.22\textwidth}
  >{\raggedright\arraybackslash}X
  >{\raggedright\arraybackslash}p{0.25\textwidth}@{}}
\toprule
\textbf{Operation} & \textbf{Primary action} &
\textbf{Fiber contracted by the output} & \textbf{Memory of reconstruction}\\
\midrule
Sum & aggregation or translation & alternatives with the same evaluation &
origin, coefficients and support\\
Product & composition or scale & nonunitary factorizations and classes &
order of factors, quotient and sheet\\
Power & iteration & composition orbits &
exponent, order, and initial point\\
Root & multivalued inversion & antecedents of a power &
branch, phase, orientation, and boundary\\
Logarithm & depth coordinate & periods and determinations &
base, generator, domain, branch, and cutoff\\
Cobordism and metric derivative & local difference and normalization & modo constante &
incidence, orientation, length, scale and reference\\
Integration & path accumulation & paths with equal published value &
path, measure, boundary and holonomy\\
\bottomrule
\end{tabularx}
\end{center}

TPK conserves these coordinates within transportable states and not as a
list added to the result. Its file \((D,C,f,I,L,M)\) declares for each
operator the domain, codomain, rule, invariants, the information
published or lost, and the conserved memory. The transition

\[
 \text{visible operation}
 \longrightarrow
 \text{lift with memory}
 \longrightarrow
 \text{compatible section}
\]

provides the link between the inaugural arithmetic, reconstruction of the
continuum, and the subsequent definition of genealogical number. The
number--particle boundary will preserve the same architecture when adding a
dimensional realization to the surviving section.

\paragraph{Inaugural domain of the nonadic cell and causal exclusion of APP, TRIT, TPK, and analytic constants}
The priority of sum and product in APP, the conservation of phase, sheet,
carry, and boundary by TPK, and the transition from section to route constitute the
premises of the lift
\eqref{eq:levantamiento-genealogico-operacion-rev2}, the notion of completeness
over fibers, and the table of memories. An operation is not identified with its physical realization:
each subsequent step preserves its domain, its chart, and its source.
